% Ad Atticum 2.5 --- headnote
% ad-atticum-02-05, mid-April 59 BC, written at Antium

\textit{Cicero to Atticus, written at Antium in mid-April 59
BC. The triumvirs (``these as the senders'') have offered
Cicero a free legateship to Egypt --- the journey he had
long wished to make, since he had not seen the East ---
and Cicero is on the ground of refusing it. The Iliadic
tags are the keys: ``I am ashamed before the Trojans and
the long-robed Trojan women'' (Hector, before going out
to face Achilles); ``Polydamas first will fasten reproach
on me'' (Hector again, of how the Trojans will judge him
if he flees). The opinion of Cato ``alone to me as a
hundred thousand'' is the present anxiety; the opinion of
``the histories six hundred years hence'' the deeper one.
Cicero closes the paragraph by noting that there is some
glory also in not accepting, and asks Atticus, if Pompey's
agent Theophanes happens to broach the subject, not to
shut it utterly down.}

\textit{The other half is the consular field for 58 BC ---
``whether Pompey and Crassus, or Servius Sulpicius with
Gabinius'' --- and the augurate, the one office in which
Cicero might be tempted by them, plus the reform of
Publius Clodius's tribunate-by-adoption. The letter ends
on the now-familiar resolve: ``I shall philosophize.''}
